% Ad Familiares 7.15 --- headnote
% ad-familiares-07-15, June 53 BC, Rome

\textit{Cicero to C. Trebatius Testa in Gaul, written from Rome
in June 53 BC. (The manuscript order of the seventh book of the
\textit{Familiares} places this letter after 7.14, but the
dateline puts it about a month earlier --- Tyrrell and Purser,
following the Loeb sequence, retain the manuscript position.)
The letter is brief and uncharacteristically tender. Trebatius
has at last written that he is enjoying himself in Gaul; Cicero,
who has spent the preceding letters of the series scolding him
for homesickness, now finds himself stung by the opposite fact.
The opening sentence is the whole joke: ``how peevish a thing it
is to be in love can be gathered from this alone.'' The mood is
banter still, but the affection beneath it is plain.}

\textit{The substance lies in the second paragraph: Trebatius
has formed a close friendship with C. Matius, the cultivated
Roman knight who appears repeatedly in Cicero's correspondence
as a man of charm and learning, and who would later be the
recipient of Cicero's defence of his own conduct after Caesar's
murder (\textit{Fam.} 11.27--28). Cicero's emphasis on the
acquisition is sincere: of all the things Trebatius might bring
home from Gaul --- gold from Britain, a cohort of \textit{essedarii},
a fat purse from Caesar --- a Matius is by far the most precious.
The closing \textit{cura ut valeas} is the formula of intimate
correspondence.}
